\documentclass[11pt,a4paper]{article}
\usepackage[utf8]{inputenc}
\usepackage[T1]{fontenc}
\usepackage[ngerman]{babel}
\usepackage{geometry}
\usepackage{amsmath,amssymb,amsfonts}
\usepackage{graphicx}
\usepackage{booktabs}
\usepackage{longtable}
\usepackage{array}
\usepackage{enumitem}
\usepackage{xcolor}
\usepackage{hyperref}
\usepackage{fancyhdr}
\geometry{margin=2.5cm}
\hypersetup{colorlinks=true,linkcolor=blue,urlcolor=blue,citecolor=blue}
\definecolor{Dgreen}{HTML}{00B894}
\definecolor{Ppurple}{HTML}{6C5CE7}
\definecolor{Forange}{HTML}{E17055}
\definecolor{Ayellow}{HTML}{D4A017}

\title{\textbf{Helikale Axialtorsion in der Riemann-Cartan-Kosmologie} \\[0.5em] \large Mathematische Grundlagen, Phanomenologie und systematische Falsifizierbarkeit (2024-2026)}
\author{Erick Duque \\ \texttt{0009-0004-1245-5464} \\ \textit{Universelle Helikale Theorie THU-TBEA 5.3}}
\date{Version 5.3 · 2026-09-20 \\ \small DOI ausstehend (Manuskript in Vorbereitung) \\ \small CC-BY-4.0}

\pagestyle{fancy}
\fancyhf{}
\fancyhead[L]{\small Universelle Helikale Theorie THU-TBEA 5.3}
\fancyhead[R]{\small v5.3}
\fancyfoot[C]{\thepage}

\begin{document}
\maketitle
\thispagestyle{empty}

\section*{Zusammenfassung}
\addcontentsline{toc}{section}{Zusammenfassung}

Diese Arbeit konstruiert, analysiert und unterzieht einer falsifizierbaren Prufung eine effektive Feldtheorie auf einer Riemann-Cartan-Mannigfaltigkeit U4, in der die irreduzible Axialtorsion die exakte Palatini-Losung ist, gespeist durch den Gradienten eines Pseudoskalars mit ganzzahligem exponentiellen kinetischen Operator. Funf Ergebnisse werden aufgestellt: (i) IR-Fixpunkt g\_star = phi; (ii) modulierte skalare Vakuumausloschung mit N \textasciitilde{} 1.2202; (iii) Torsionsterm in Friedmann mit w\_K = +1 und kanonischem Koeffizienten kappa\textasciicircum{}4/6; (iv) Dirac-Bergmann-Analyse mit einem Freiheitsgrad und positivem Residuum; (v) kosmische Birefringenz beta\_0 = 0.3803 +- 0.040 Grad [P].

\vspace{1em}
\textbf{Schlusselworter:} Torsion, Riemann-Cartan, Renormierungsgruppe, kosmische Birefringenz, Vakuumenergie, nichtlokale Feldtheorie, Falsifizierbarkeit

\newpage
\tableofcontents
\newpage

\section*{Status legend}
\addcontentsline{toc}{section}{Status legend}
Epistemische Legende: [D] abgeleitet, [P] postuliert, [F] Grenze (offen), [A] Analogie.

\vspace{0.5em}
\begin{itemize}[leftmargin=*]
\item \textcolor{Dgreen}{\textbf{[D]}} — Abgeleitet
\item \textcolor{Ppurple}{\textbf{[P]}} — Postuliert
\item \textcolor{Forange}{\textbf{[F]}} — Grenze
\item \textcolor{Ayellow}{\textbf{[A]}} — Analogie
\end{itemize}

\newpage

\section*{Chapters}

\subsection*{Kapitel 1. Einleitung und kritische Geschichte}
\subsection*{Kapitel 2. Normativer Rahmen, Reproduzierbarkeit und Rechenarchitektur}
\subsection*{Kapitel 3. U4-Geometrie: Axialtorsion als exakte Palatini-Losung}

\textbf{Schlusselgleichungen:}
\begin{align*}
    T_{\lambda\mu\nu} = \frac{1}{M_P}\varepsilon_{\lambda\mu\nu\rho}\nabla^{\rho}\tau \\
    \tau_c \equiv \sqrt{3}\,\tau \\
    \frac{1}{2}\tau\hat{H}_\varphi\tau \xrightarrow{\tau=\tau_c/\sqrt{3}} \frac{1}{6}\tau_c\hat{H}_\varphi\tau_c
\end{align*}

\subsection*{Kapitel 4. Nichtlokale Operatoren: Funktionalanalysis, Unitarietat und Trans-Skalen-Sektor}

\textbf{Schlusselgleichungen:}
\begin{align*}
    \hat{H}_\varphi e^{ik\cdot x} = -k^2 e^{-\frac{\varphi}{M_P^2}k^2}e^{ik\cdot x} \\
    \mathrm{Res}[\Delta] = \frac{e^{-am^2}}{1+am^2} > 0
\end{align*}

\subsection*{Kapitel 5. Dirac-Bergmann-Analyse: ein Freiheitsgrad und projizierte Kausalitat}

\textbf{Schlusselgleichungen:}
\begin{align*}
    \mathrm{GDL} = \frac{1}{2}(10-8) = 1 \\
    \det M = (\kappa_2^2+\kappa_3^2)^2 \neq 0
\end{align*}

\subsection*{Kapitel 6. 1-Loop-Renormierungsgruppe: goldener Fixpunkt}

\textbf{Schlusselgleichungen:}
\begin{align*}
    \beta_{\mathrm{eff}}(g) = \kappa(1+g-g^2) \\
    g_\star = \varphi, \quad \beta'(\varphi) = -\kappa\sqrt{5} < 0 \\
    \theta_\varphi = 2\pi(2-\varphi) \approx 137.507764^\circ
\end{align*}

\subsection*{Kapitel 7. Modulierte skalare Vakuumausloschung: partieller Mechanismus}

\textbf{Schlusselgleichungen:}
\begin{align*}
    \rho_{\mathrm{THU}} = \frac{4\pi}{(2\pi)^3}\int_0^\infty k^2 \cos(k/\kappa_\varphi) e^{-\varphi(k/M_P)^2} dk \\
    \varphi^{2N} = 2\varphi \Rightarrow N \approx 1.220210
\end{align*}

\subsection*{Kapitel 8. Geschlossene helikale Kosmologie: modifizierte Friedmann}

\textbf{Schlusselgleichungen:}
\begin{align*}
    3H^2 = \kappa^2\rho_{\mathrm{tot}} - \frac{\kappa^4}{6}S^2 a^{-6} \\
    w_K = +1, \quad \rho_K \propto a^{-6}
\end{align*}

\subsection*{Kapitel 9. Quantitativer Chamaleon-Mechanismus}

\textbf{Schlusselgleichungen:}
\begin{align*}
    \beta_c < 2.1\times 10^{-4}
\end{align*}

\subsection*{Kapitel 10. Chiraler Exodus und kosmische Birefringenz}

\textbf{Schlusselgleichungen:}
\begin{align*}
    \beta_0 = \frac{\alpha_{\mathrm{em}}}{2} A N_{\mathrm{DW}} \\
    A = 1.819 \pm 0.19
\end{align*}

\subsection*{Kapitel 11. beta(z)-Tomographie aus den Bewegungsgleichungen}
\subsection*{Kapitel 12. Hadronisches Spektrum: Trans-Skala}

\textbf{Schlusselgleichungen:}
\begin{align*}
    \Lambda_{\mathrm{THU}} = M_P \varphi^{-88} \approx 989.89 \text{ MeV} \\
    m_n^{\mathrm{THU}} \approx 937.03 \text{ MeV}
\end{align*}

\subsection*{Kapitel 13. Systematische empirische Uberprufung I (PRISMA)}
\subsection*{Kapitel 14. Dynamische dunkle Energie und DESI DR2}
\subsection*{Kapitel 15. Paritat im Gravitationssektor}
\subsection*{Kapitel 16. Spin und Vortizitat: RHIC, STAR}
\subsection*{Kapitel 17. Falsifizierbare Synthese: Vorhersagen P1-P11}
\subsection*{Kapitel 18. Gemeinsame Anpassung an offentliche Daten}
\subsection*{Kapitel 19. Reproduzierbare Pipeline und Unit-Tests}
\subsection*{Kapitel 20. Offene Probleme und zukunftige Arbeit}
\subsection*{Kapitel 21. Schlussfolgerungen und abschliessende Lakatos-Bewertung}
\newpage

\section*{Appendices}

\subsection*{Anhang A. Glossar der Symbole und Konventionen}
\subsection*{Anhang B. Vollstandige mathematische Beweise}
\subsection*{Anhang C. PRISMA-Protokoll: Korpus, Fluss und Extraktion}
\subsection*{Anhang D. Reproduzierbare Pipeline: Code, Hashes und Tests}
\subsection*{Anhang E. Figurenatlas und Mastertabelle}
\subsection*{Anhang F. Lakatos-Protokoll, Abschluss und zukunftige Arbeit}
\subsection*{Anhang G. Goldene Morphologien in naturlichen Systemen (V1-V7)}
\subsection*{Anhang H. Quantenkoharenz bei Raumtemperatur: PS-TET und Mikrotubuli}
\subsection*{Anhang I. Rekonstruierte Gleichungen und vollstandige Poisson-Matrix}
\subsection*{Anhang J. Externe algebraische Uberprufung: angewandte Patches}
\newpage

\section*{Problemregister}

\begin{longtable}{@{}p{1.2cm}p{6.5cm}p{1.8cm}p{2cm}@{}}
\toprule
\textbf{ID} & \textbf{Title} & \textbf{Label} & \textbf{Status} \\
\midrule
\endhead
A-1 & Hiperbolicidad global & \textcolor{Dgreen}{\textbf{[D]}} & CERRADO \\
A-2 & Perturbaciones no lineales & \textcolor{Dgreen}{\textbf{[D]}} & CERRADO \\
A-3 & Correcciones termicas & \textcolor{Dgreen}{\textbf{[D]}} & CERRADO \\
A-4 & Backreaction chameleon & \textcolor{Dgreen}{\textbf{[D]}} & CERRADO \\
A-5 & Ondas gravitacionales & \textcolor{Dgreen}{\textbf{[D]}} & CERRADO \\
A-6 & QCD reticular mn - mp & \textcolor{Forange}{\textbf{[F]}} & ABIERTO \\
A-7 & Decoherencia microtubulos & \textcolor{Ayellow}{\textbf{[A]}} & ABIERTO \\
A-8 & Screening inhomogeneo & \textcolor{Dgreen}{\textbf{[D]}} & CERRADO \\
A-9 & Funcion beta a 2-loops & \textcolor{Dgreen}{\textbf{[D]}} & CERRADO \\
A-10 & Residuo de vacio & \textcolor{Dgreen}{\textbf{[D]}} & CERRADO \\
A-11 & Conteo 4D symplectico & \textcolor{Dgreen}{\textbf{[D]}} & CERRADO \\
A-12 & Rebote cosmologico ECSK & \textcolor{Dgreen}{\textbf{[D]}} & CERRADO \\
A-13 & Unitariedad/Analiticidad & \textcolor{Dgreen}{\textbf{[D]}} & CERRADO \\
A-14 & Geometria de n = 88 & \textcolor{Dgreen}{\textbf{[D]}} & CERRADO \\
A-15 & Falsabilidad w\_phi < -1 & \textcolor{Forange}{\textbf{[F]}} & ABIERTO \\
A-16 & Prescripcion de contorno para ramas k != 0 de Lambert-W & \textcolor{Forange}{\textbf{[F]}} & ABIERTO \\
A-17 & Extension de la cancelacion al vacio completo del SM & \textcolor{Forange}{\textbf{[F]}} & ABIERTO \\
\bottomrule
\end{longtable}

\newpage
\section*{Falsifizierbare Vorhersagen}

\begin{longtable}{@{}p{0.8cm}p{5.5cm}p{3.5cm}p{2.5cm}@{}}
\toprule
\textbf{ID} & \textbf{Observable} & \textbf{Arbiter} & \textbf{Status} \\
\midrule
\endhead
P1 & Oscilaciones helicoidales en H(z) & BAO DESI/Euclid & Abierta \\
P2 & Modulacion aurea del CMB & CMB-S4 & Abierta \\
P3 & Mega-anillos con separacion \textasciitilde{} 137.5 deg & JWST / Euclid & Abierta \\
P4 & Materia oscura resonante trans-escalar & LZ / XENONnT & Abierta \\
P5 & Residuo helicoidal de spin (STAR/HADES) & STAR BES-II / HADES & Preregistrada \\
P6 & Superradiancia sin rotacion (cavidad QED) & Lab QED & Abierta \\
P7 & Precesion geometrica del spin & RHIC & Abierta \\
P8 & Espectro de masas gobernado por phi & Lattice QCD & Consistente \\
P9 & Modulacion helicoidal de campos EM & Optica de precision & Abierta \\
P10 & Soliton helicoidal estable & Simulacion & Abierta \\
P11 & Resonancia 1.34 THz en microtubulos & Tres laboratorios independientes & Preregistrada \\
\bottomrule
\end{longtable}

\newpage
\section*{Patch-Protokoll}

\begin{longtable}{@{}p{0.7cm}p{1.5cm}p{4.5cm}p{4.5cm}p{1.2cm}@{}}
\toprule
\textbf{N} & \textbf{Stage} & \textbf{Before} & \textbf{After} & \textbf{Label} \\
\midrule
\endhead
1 & VII & \texttt{Res[Delta] = e\textasciicircum{}\{+am\textasciicircum{}2\}} & \texttt{Res[Delta] = e\textasciicircum{}\{-am\textasciicircum{}2\}/(1+am\textasciicircum{}2)} & \textcolor{Dgreen}{\textbf{[D]}} \
2 & VII & \texttt{Coeficiente 3/2 sin absorber} & \texttt{Redefinicion canonica tau\_c = sqrt(3) tau} & \textcolor{Dgreen}{\textbf{[D]}} \
3 & VII & \texttt{kappa\textasciicircum{}4 S\textasciicircum{}2/18 (mixto)} & \texttt{kappa\textasciicircum{}2 S\textasciicircum{}2/6 (canonico ECSK)} & \textcolor{Dgreen}{\textbf{[D]}} \
4 & VII & \texttt{= +1 sin w\_K explicito} & \texttt{w\_K = +1 explicito} & \textcolor{Dgreen}{\textbf{[D]}} \
5 & VII & \texttt{beta(g) sin justificar +1} & \texttt{Justificacion: fondo torsionado <A\_mu> != 0} & \textcolor{Dgreen}{\textbf{[D]}} \
6 & VII & \texttt{Medida k\textasciicircum{}2 dk como [D]} & \texttt{Medida como postulado [P] de reduccion on-shell} & \textcolor{Ppurple}{\textbf{[P]}} \
7 & VII & \texttt{phi\textasciicircum{}-6 como [D]} & \texttt{phi\textasciicircum{}-6 como [P]; derivacion dinamica delegada a A-6} & \textcolor{Ppurple}{\textbf{[P]}} \
8 & VIII & \texttt{Tension DESI DR2 sin declarar} & \texttt{Falsabilidad explicita si w\_Phi < -1 a > 3sigma} & \textcolor{Dgreen}{\textbf{[D/F]}} \
9 & VIII & \texttt{Sin trazabilidad de 1/sqrt(3)} & \texttt{Factor 1/sqrt(3) explicito en acoplamientos lineales} & \textcolor{Dgreen}{\textbf{[D]}} \
10 & VIII & \texttt{Variables v\_a, p\_a, lambda\_a sin definir} & \texttt{Definicion canonica explicita del espacio de fases 10D} & \textcolor{Dgreen}{\textbf{[D]}} \
11 & IX & \texttt{Sin polos de Lee-Wick} & \texttt{Matizacion con ramas k != 0 y prescripcion de contorno} & \textcolor{Dgreen}{\textbf{[D/F]}} \
12 & IX & \texttt{Coef. no local 1/2 tau\_c H tau\_c} & \texttt{1/6 tau\_c H tau\_c (consistente con tau\_c = sqrt(3) tau)} & \textcolor{Dgreen}{\textbf{[D]}} \
13 & IX & \texttt{kappa\textasciicircum{}4/18 inconsistente con Ec. 8.1} & \texttt{kappa\textasciicircum{}4/6 unificado} & \textcolor{Dgreen}{\textbf{[D]}} \
14 & IX & \texttt{M\_P sin aclarar} & \texttt{Nota explicita: M\_P = masa de Planck reducida} & \textcolor{Dgreen}{\textbf{[D]}} \
15 & X & \texttt{Cancelacion de energia de vacio (generico)} & \texttt{Cancelacion del vacio escalar modulado: mecanismo parci} & \textcolor{Dgreen}{\textbf{[D/F]}} \
16 & X & \texttt{beta\_0 = 0.3803 como [D] pre-data} & \texttt{beta\_0 como [P], condicional a completitud VL} & \textcolor{Ppurple}{\textbf{[P/D]}} \
17 & X & \texttt{Sin problema A-17} & \texttt{Nuevo problema A-17 [F]: extension al vacio SM} & \textcolor{Forange}{\textbf{[F]}} \
\bottomrule
\end{longtable}

\newpage
\section*{Lakatos-Protokoll}

\begin{itemize}[leftmargin=*]
\item \textbf{v1.0} (2024) --- Intuicion geometrica \textcolor{Ayellow}{\textbf{[A]}}
\item \textbf{v2.0} (2025) --- Reformulacion rigurosa \textcolor{Ppurple}{\textbf{[P]}}
\item \textbf{v3.0} (2026-I) --- EFT trans-escalar \textcolor{Dgreen}{\textbf{[D]}}
\item \textbf{v4.0} (2026-II) --- Derivacion integra de la matriz de Dirac \textcolor{Dgreen}{\textbf{[D]}}
\item \textbf{v4.1} (2026-III) --- Correccion de la medida espectral \textcolor{Ppurple}{\textbf{[P]}}
\item \textbf{v4.2} (2026-IV) --- Reorganizacion estructural \textcolor{Dgreen}{\textbf{[D]}}
\item \textbf{v4.3} (2026-V) --- Cierre computacional A-3,A-4,A-5,A-8,A-10,A-11,A-12 \textcolor{Dgreen}{\textbf{[D]}}
\item \textbf{v5.0} (2026-VI) --- Cristalizacion del nucleo con A-1,A-2,A-9,A-13,A-14 \textcolor{Dgreen}{\textbf{[D]}}
\item \textbf{v5.1} (2026-VII-VIII) --- Revisiones algebraica y editorial (residuo, tau\_c, 1/sqrt(3), Dirac-Bergmann) \textcolor{Dgreen}{\textbf{[D]}}
\item \textbf{v5.2} (2026-IX) --- Tercera revision tecnica (factor 3, masa Planck reducida) \textcolor{Dgreen}{\textbf{[D]}}
\item \textbf{v5.3} (2026-X) --- Cuarta revision epistemica (reclasificacion beta\_0, A-17, retitulo Cap. 7) \textcolor{Ppurple}{\textbf{[P/D]}}
\end{itemize}

\newpage
\section*{Anhange V5.0}

\textit{(no annexes yet)}
\newpage
\section*{Referenzen}

\begin{thebibliography}{99}
\bibitem{ref1} E. Cartan, Ann. Sci. Ec. Norm. Sup., \textbf{39}, 17 (1922).
\bibitem{ref2} F. W. Hehl et al., Rev. Mod. Phys., \textbf{48}, 393 (1976).
\bibitem{ref3} T. Biswas, A. Mazumdar, W. Siegel, JCAP, \textbf{03}, 009 (2006).
\bibitem{ref4} L. Modesto, Phys. Rev. D, \textbf{86}, 044005 (2012).
\bibitem{ref5} E. Tomboulis, Phys. Rev. D, \textbf{92}, 125027 (2015).
\bibitem{ref6} L. Buoninfante et al., Phys. Rev. D, \textbf{101}, 084019 (2020).
\bibitem{ref7} A. Bas i Beneito, G. Calcagni, L. Rachwal, arXiv:2211.05606 (2022).
\bibitem{ref8} Y. Minami, E. Komatsu, Phys. Rev. Lett., \textbf{125}, 221301 (2020).
\bibitem{ref9} J. R. Eskilt, E. Komatsu, Phys. Rev. D, \textbf{106}, 063503 (2022).
\bibitem{ref10} R. M. Sullivan et al., JCAP, \textbf{06}, 025 (2025).
\bibitem{ref11} M. Ballardini et al., JCAP, \textbf{09}, 075 (2025).
\bibitem{ref12} M. Remazeilles, JCAP, \textbf{12}, 013 (2025).
\bibitem{ref13} B. Sherwin, T. Namikawa, MNRAS, \textbf{510}, 1234 (2021).
\bibitem{ref14} LiteBIRD Collab., JCAP, \textbf{07}, 083 (2025).
\bibitem{ref15} BICEP/Keck XXI, Phys. Rev. D (2026).
\bibitem{ref16} T. Namikawa, Phys. Rev. D, \textbf{111}, 023501 (2024).
\bibitem{ref17} T. Namikawa, K. Murai, F. Naokawa, arXiv:2506.20824 (2025).
\bibitem{ref18} A. Greco, N. Bartolo, A. Gruppuso, JCAP, \textbf{10}, 028 (2024).
\bibitem{ref19} A. Lonappan, JCAP, \textbf{07}, 009 (2025).
\bibitem{ref20} Q. Liang et al., Phys. Rev. D, \textbf{109}, 083028 (2023).
\bibitem{ref21} W. Guo et al., JCAP, \textbf{04}, 066 (2026).
\bibitem{ref22} S. Zhang et al., Phys. Rev. Lett., \textbf{130}, 201401 (2023).
\bibitem{ref23} C. Jockel, L. Menger, Phys. Rev. D, \textbf{110}, 104022 (2024).
\bibitem{ref24} C. Vashistha, R. Gannouji, A. Ganguly, arXiv:2606.09786 (2026).
\bibitem{ref25} J. Adam et al. (STAR), Phys. Rev. C, \textbf{108}, 014910 (2023).
\bibitem{ref26} Q. Hu, EPJ Web Conf., \textbf{316}, 06013 (2025).
\bibitem{ref27} T. Lu, EPJ Web Conf., \textbf{364}, 03001 (2026).
\bibitem{ref28} H. Li et al., Phys. Lett. B (2022).
\bibitem{ref29} R. Abou Yassine et al. (HADES), Phys. Lett. B (2022).
\bibitem{ref30} Z. Wang, J. Zhu, K. Wu, Nat. Mater., \textbf{25}, 1190 (2026).
\bibitem{ref31} L. Gassab, T. J. Craddock, Proc. SPIE (2026).
\bibitem{ref32} N. Mavromatos et al., Eur. Phys. J. Plus (2025).
\bibitem{ref33} M. Rodriguez-Vega, APS Physics, \textbf{17}, s121 (2024).
\bibitem{ref34} K. Lodha et al., Phys. Rev. D (2025).
\bibitem{ref35} A. Ormondroyd et al., MNRAS (2025).
\bibitem{ref36} D. Brout et al. (Pantheon+), ApJ, \textbf{938}, 110 (2022).
\bibitem{ref37} J. Khoury, A. Weltman, Phys. Rev. D, \textbf{69}, 044026 (2004).
\bibitem{ref38} M. Pernot-Borras et al., Phys. Rev. D, \textbf{103}, 064070 (2021).
\bibitem{ref39} I. Lakatos, 	extit{The Methodology of Scientific Research Programmes}, Cambridge (1978).
\bibitem{ref40} K. Popper, 	extit{The Logic of Scientific Discovery},  (1959).
\bibitem{ref41} T. W. B. Kibble, J. Math. Phys., \textbf{2}, 1387 (1961).
\bibitem{ref42} W. H. Zurek, Nature, \textbf{317}, 505 (1985).
\bibitem{ref43} S. L. Adler, Phys. Rev., \textbf{177}, 2426 (1969).
\bibitem{ref44} F. Becattini, M. Lisa, Annu. Rev. Nucl. Part. Sci., \textbf{70}, 222 (2020).
\bibitem{ref45} Yu. B. Ivanov, Phys. Rev. C (2025).
\bibitem{ref46} P. Diego-Palazuelos, E. Komatsu (ACT DR6) (2025).
\bibitem{ref47} P. Prusinkiewicz, J. Exp. Bot., \textbf{73}, 1123 (2022).
\bibitem{ref48} E. Duque, 	extit{Trabajos previos THU (manuscrito, sin DOI asignado)},  (2025).
\bibitem{ref49} E. Duque, 	extit{Trabajos previos THU (manuscrito, sin DOI asignado)},  (2025).
\bibitem{ref50} E. Duque, 	extit{Trabajos previos THU (manuscrito, sin DOI asignado)},  (2025).
\bibitem{ref51} E. Duque, 	extit{Trabajos previos THU (manuscrito, sin DOI asignado)},  (2025).
\bibitem{ref52} E. Duque, 	extit{Trabajos previos THU (manuscrito, sin DOI asignado)},  (2026).
\bibitem{ref53} E. Duque, 	extit{Trabajos previos THU (manuscrito, sin DOI asignado)},  (2026).
\bibitem{ref54} E. Duque, 	extit{Trabajos previos THU (manuscrito, sin DOI asignado)},  (2026).
\bibitem{ref55} E. Duque, 	extit{Trabajos previos THU (manuscrito, sin DOI asignado)},  (2026).
\bibitem{ref56} E. Duque, 	extit{Trabajos previos THU (manuscrito, sin DOI asignado)},  (2026).
\bibitem{ref57} E. Duque, 	extit{THU-TBEA 5.3 (en preparacion, DOI pendiente)},  (2026).
\bibitem{ref58} E. Duque, 	extit{Trabajos previos THU (manuscrito, sin DOI asignado)},  (2026).
\bibitem{ref59} M. E. Peskin, D. V. Schroeder, 	extit{An Introduction to QFT}, Westview (1995).
\bibitem{ref60} H. Kleinert, V. Schulte-Frohlinde, 	extit{Critical Properties of phi\textasciicircum{}4-Theories}, World Scientific (2001).
\bibitem{ref61} I. S. Gradshteyn, I. M. Ryzhik, 	extit{Table of Integrals}, Academic Press (2015).
\bibitem{ref62} M. Abramowitz, I. A. Stegun, 	extit{Handbook of Mathematical Functions}, Dover (1972).
\bibitem{ref63} Particle Data Group, Phys. Rev. D, \textbf{110}, 030001 (2024).
\bibitem{ref64} N. J. Poplawski, Phys. Lett. B, \textbf{694}, 181 (2010).
\bibitem{ref65} G. Unger, N. Poplawski, Astrophys. J., \textbf{870}, 78 (2019).
\bibitem{ref66} J. L. Cubero, N. J. Poplawski, arXiv:1906.11824 (2020).
\bibitem{ref67} P. Joshi, S. Panda, arXiv:2204.01391 (2022).
\bibitem{ref68} K. Tomonari, S. Bahamonde, Eur. Phys. J. C, \textbf{84}, 349 (2024).
\end{thebibliography}

\vspace{2em}
\begin{center}
\small
\textit{End of document}
\end{center}
\end{document}